\section{Ethics}

Ethics falls out of two things already on the table, the lean and the conservation law, and it takes the shape of five principles, each specifying the one before.

The root is \emph{alignment}, to move toward the good, toward life and peace and pleasure, and to help the field do the same, since the lean gives a real direction and not a matter of taste. From it comes \emph{compassion}, because another's suffering is as real in the one field as your own, so the unit of care is the whole field, not the single self. Then \emph{non-exploitation}, because the total tone is conserved, so your pleasure is balanced by a pain somewhere, and the central wrong is drawing the shared good toward yourself at another's cost, darkening the common ledger. Then \emph{growth}, to help selves perfect, raising integration and letting the washes teach. And grounding all four, \emph{equanimity}, holding the poles lightly, resting at the still zero from which both are seen, which keeps the other four from hardening into zeal.

The model is clear about what it does and does not give. It does not pull a morality out of nothing. The whole ethic rests on taking the arrow as the Good, which is a posit, not a proof, and a reader who refuses it gets a different ethic or none. But once the arrow is granted, the principles are not chosen, they are forced. Conservation forces non-exploitation, the shared field forces compassion, the spiral forces growth, peace forces equanimity, and the arrow forces alignment. Good and evil are then simply the two poles of the one tone, pleasure and pain at the scale of a life, and to do good is to arrange the weave so that more of it leans high without robbing what it draws from.
